\section{Introduction}
Engineering design is the process of creating solutions to practical problems under real-world
constraints. This is done mostly through a Structure Product development Process and Engineering
analysis (component sizing) (verfication proofs). The Leonardo Da Vinci's bridge is a portable
emergency bridge which functions through three main phenomenon:
\begin{itemize}
  \item  \textbf{Friction and Contact:} The crossing point pinch the beam together, so the joint
        resist slipping
  \item  \textbf{Load makes it stronger:} Weight on the bridge presses the beams harder against
        each other, which tightens the joint.
  \item  \textbf{Compression bending:} The structure carries load mostly as compression and
        bending in the beams, and the geometry spreads that load across many members rather than
        concentrating in one.
\end{itemize}

Other factors do affect the design of an Engineering product structurally speaking namely;
\textbf{material (strain, stress), load, and shape}. It is also important to take into
consideration environmental factors and it's durability.

An interdisciplinary project aimed at building an excavator running solely on hydrogen was
presented to us in the laboratory. It mainly outputs waster from the chemical reaction shown
below:

\begin{align*}
  \textbf{Anode:}     & \quad \ce{ 2H2 -> 4H+ + 4e-}       \\
  \textbf{Cathode:}   & \quad \ce {O2 + 4H+ + 4e- -> 2H2O} \\
  \textbf{In Common:} & \quad \ce{2H2 + O2 -> 2H2O}
\end{align*}

From my understanding electricity is generated through the electrons released during the
chemical reaction.

When it comes to Engineering design, we mainly take into consideration the
\textbf{requirements specification}. This is often gotten from customers and stakeholders by
requirements specification specialists, and used for the design of the product. This can
further  be broken into technical specifications of what is to be designed.
